\documentclass[11pt]{article}
\usepackage[letterpaper,margin=1in]{geometry}
\usepackage[T1]{fontenc}
\usepackage[utf8]{inputenc}
\usepackage{lmodern,microtype}
\usepackage{amsmath,amssymb,mathtools,bm,slashed}
\usepackage{xcolor,enumitem,booktabs,array,needspace}
\usepackage[most]{tcolorbox}
\usepackage{fancyhdr,titlesec,parskip,hyperref}
\definecolor{qftblue}{HTML}{245A8D}
\definecolor{qftgreen}{HTML}{2F6B4F}
\definecolor{qftgray}{HTML}{F4F6F8}
\hypersetup{colorlinks=true,linkcolor=qftblue,urlcolor=qftblue,pdfauthor={Nima Lashkari}}
\pagestyle{fancy}\fancyhf{}
\lhead{PHYS 662: Quantum Field Theory I}\cfoot{\thepage}
\setlength{\headheight}{14pt}
\titleformat{\section}{\Large\bfseries\color{qftblue}}{\thesection.}{0.6em}{}
\titleformat{\subsection}{\large\bfseries\color{qftgreen}}{\thesubsection}{0.6em}{}
\setlist[enumerate]{leftmargin=2em,itemsep=0.25em,topsep=0.35em}
\setlist[itemize]{leftmargin=2em,itemsep=0.2em,topsep=0.35em}
\newcommand{\dd}{\mathrm d}
\newcommand{\ii}{\mathrm i}
\newcommand{\eps}{\varepsilon}
\newcommand{\tr}{\operatorname{tr}}
\newcommand{\Tr}{\operatorname{Tr}}
\newcommand{\diag}{\operatorname{diag}}
\newcommand{\Lag}{\mathcal L}
\newcommand{\D}{\mathcal D}
\newcommand{\vac}{\lvert0\rangle}
\newcommand{\vbra}{\langle0\rvert}
\newcommand{\Tord}{\mathrm T}
\newcommand{\ket}[1]{\lvert#1\rangle}
\newcommand{\bra}[1]{\langle#1\rvert}
\newtcolorbox{keyidea}[1][]{enhanced,breakable,colback=qftgray,colframe=qftblue,
 boxrule=0.7pt,arc=1.5mm,title={#1},fonttitle=\bfseries,coltitle=white}
\newtcolorbox{remarkbox}[1][]{enhanced,breakable,colback=white,colframe=qftgreen,
 boxrule=0.6pt,arc=1.5mm,title={#1},fonttitle=\bfseries,coltitle=white}
\newcommand{\intp}{\int\!\frac{\dd^4p}{(2\pi)^4}}
\newcommand{\MSbar}{\overline{\mathrm{MS}}}
\newcommand{\module}[1]{\clearpage\section{#1}}
\newcommand{\braket}[2]{\langle#1\mid#2\rangle}
\newcommand{\dPi}{\dd\Pi}
\newcommand{\cM}{\mathcal M}
\newcommand{\cH}{\mathcal H}
\newcommand{\cS}{\mathcal S}
\newcommand{\Zp}{Z_{\rm pole}}
\newcommand{\Kallen}{\mathcal K}
\newcommand{\lecturefront}[4]{%
\rhead{Lecture #1}\hypersetup{pdftitle={PHYS 662 Lecture #1: #2}}
\begin{document}\pagenumbering{gobble}
\begin{titlepage}\centering\vspace*{1.05in}
{\scshape\Large PHYS 662 --- Quantum Field Theory I\par}
\vspace{0.35in}{\Huge\bfseries\color{qftblue}Lecture #1\par}
\vspace{0.2in}{\LARGE\bfseries #2\par}
\vspace{0.2in}{\large\color{qftgreen}#3\par}
\vspace{0.6in}{\large Nima Lashkari\par}{\large Fall 2026\par}
\vfill\begin{minipage}{0.86\textwidth}\small #4\par\medskip
Course page: \url{https://nimalashkari.com/quantum-field-theory-2026-part-i/}
\end{minipage}\vspace{0.5in}\end{titlepage}
\pagenumbering{roman}\setcounter{tocdepth}{1}\tableofcontents
\vfill\begin{remarkbox}[Place in the course]
Lectures 21--22 developed interacting correlators and scalar loops.
Lecture 23 defines the scattering states and the S-matrix.
Lecture 24 converts amplitudes into rates and cross sections.
Lecture 25 derives the reduction formula connecting these subjects.
Homework 8 covers this three-lecture sequence.
\end{remarkbox}
\small\textbf{Reading and conventions.} These notes use the course slides and
D.~Harlow, \href{https://www.mit.edu/~harlow/HarlowQFT1.pdf}{\textit{QFT I}},
Sections 8--9, with our covariant state normalization and mostly-minus metric.
Useful online companions are D.~Tong,
\href{https://www.damtp.cam.ac.uk/user/tong/qft/qfthtml/S3.html}{\textit{QFT}, Sections 3.6--3.7},
and D.~Buchholz and W.~Dybalski,
\href{https://arxiv.org/abs/math-ph/0509047}{\textit{Scattering in relativistic quantum field theory}},
Sections 1--3, for the scope of the asymptotic construction.
For textbook study, see Weinberg, \textit{The Quantum Theory of Fields},
Vol.~I, Chapter 3, and Peskin--Schroeder, \textit{An Introduction to QFT},
Chapters 4 and 7. The accompanying editorial memo records the online
sources consulted and the corrections to the slides.
\normalsize\clearpage\pagenumbering{arabic}}
\lecturefront{25}{Scattering Matrices from Correlation Functions}
{LSZ reduction, pole residues, and external particle states}
{Path integrals compute local correlation functions, whereas experiments
prepare and detect asymptotic particles. We derive the bridge between
them using wave packets and one-particle poles, keep the normalization
explicit, and test the result on scalar and fermion amplitudes.}

\section{The question answered by reduction}
In Lectures 21--22 we computed normalized time-ordered functions
\begin{equation}
 G_N(x_1,\ldots,x_N)
 =\bra{\Omega}\Tord\phi(x_1)\cdots\phi(x_N)\ket{\Omega}.
\end{equation}
Their external momenta may be off shell. A scattering amplitude instead
has physical incoming and outgoing particles with positive energies and
\(p^2=M^2\). The Lehmann--Symanzik--Zimmermann (LSZ) reduction formula
extracts the part of a correlator that propagates these particles to
arbitrarily early or late times.

Assume a stable scalar mass shell, an interpolating field with zero vacuum
expectation value, and a nonzero overlap
\begin{equation}
 \bra{\Omega}\phi(0)\ket{p}=\sqrt{\Zp},\qquad
 \widetilde G_2(p)\sim\frac{\ii\Zp}{p^2-M^2+\ii0},\qquad
 \langle p|q\rangle=2E_p(2\pi)^3\delta^{(3)}(\bm p-\bm q).
 \label{eq:pole}
\end{equation}
These are exact interacting quantities. \(M\) is the physical pole mass,
and \(\Zp\) is the residue for the field actually used in the correlator.
If that field has no overlap with the particle, it cannot reduce that
external particle, even if a different operator could.

Why do poles help? A field insertion very far from the collision can
connect to it through a stable one-particle intermediate state. The long
time integral of that state's phase produces an energy denominator.
Doing this for every external insertion isolates the simultaneous
one-particle poles. Contributions regular at an external mass shell have
no such long-lived external particle and disappear in that pole extraction.

\begin{keyidea}[The expected answer, with our factors of \(\ii\)]
For each scalar external leg, multiply the Fourier-transformed connected
correlator by
\[
 \frac{p^2-M^2}{\ii\sqrt{\Zp}}
\]
and extract its on-shell pole residue. The result is the connected
S-matrix element \(\ii(2\pi)^4\delta^{(4)}(P_f-P_i)\cM\).
The following sections derive this rule and specify its meaning.
\end{keyidea}
The on-shell instruction is not ``put \(p^2=M^2\) before canceling the
pole.'' It is a residue limit with the Feynman boundary prescription,
understood after smearing in external momenta. Unstable resonances and
infrared-sensitive charged states require the qualifications discussed
in the preceding lectures.

\module{Particle projection and the Klein--Gordon product}
Let \(f\) be a smooth positive-frequency solution of the free equation with
\emph{physical} mass \(M\):
\begin{equation}
 f(x)=\int\dPi_p f(\bm p)e^{-\ii p\cdot x},\qquad
 (\Box+M^2)f=0,\qquad
 \dPi_p=\frac{\dd^3p}{(2\pi)^3\,2E_p}.
\end{equation}
The Klein--Gordon product on such solutions is
\begin{equation}
 (f,g)_{\rm KG}=\ii\int_{x^0=t}\dd^3x\,
 f^*\overleftrightarrow{\partial_t}g
 =\int\dPi_p f(\bm p)^*g(\bm p),
 \label{eq:KGproduct}
\end{equation}
where \(u\overleftrightarrow\partial v=u\partial v-(\partial u)v\).
It is independent of \(t\). Substituting the free mode expansion shows
that this product extracts the annihilation operator of the packet.

For the interacting Heisenberg field define the finite-time projection
\begin{equation}
 a_f(t)=\frac{\ii}{\sqrt{\Zp}}\int\dd^3x\,
 f^*(x)\overleftrightarrow{\partial_t}\phi(x).
 \label{eq:projection}
\end{equation}
It is not generally time independent. Differentiate, use the equation
for \(f\), and integrate its spatial derivatives by parts. Packet decay
removes the spatial boundary term, leaving the exact finite-time identity
\begin{equation}
 \frac{\dd a_f(t)}{\dd t}
 =\frac{\ii}{\sqrt{\Zp}}\int\dd^3x\,f^*(x)(\Box+M^2)\phi(x).
\end{equation}
With the asymptotic limits understood on scattering-state matrix elements,
\begin{equation}
 \boxed{\quad a_{\rm out}(f)-a_{\rm in}(f)
 =\frac{\ii}{\sqrt{\Zp}}\int\dd^4x\,
 f^*(x)(\Box+M^2)\phi(x).\quad}
 \label{eq:adifference}
\end{equation}
The time integral carries the same smearing and limiting prescription as
the boundary expression from which it was derived.

\begin{remarkbox}[A limit of matrix elements, not a pointwise free field]
The assertion \(\phi\sim\sqrt{\Zp}\phi_{\rm in/out}\) is an asymptotic
statement after appropriate smearing and mass-shell isolation. It is not
pointwise operator convergence. A rigorous construction uses energy
filtering and, where appropriate, time averaging. The unaveraged LSZ
projection above is used in its weak, scattering-matrix-element sense;
we do not claim norm convergence of arbitrary unsmeared field operators.
\end{remarkbox}

\module{Reducing one particle and understanding the boundary terms}
Take an outgoing packet \(f\) in addition to the outgoing configuration
\(\beta\). Equation \eqref{eq:adifference} gives
\begin{align}
 \langle\beta,\mathrm{out}|a_{\rm out}(f)|\alpha,\mathrm{in}\rangle
 &=\langle\beta,\mathrm{out}|a_{\rm in}(f)|\alpha,\mathrm{in}\rangle\notag\\
 &\quad+\frac{\ii}{\sqrt{\Zp}}\int\dd^4x\,f^*(x)
 (\Box_x+M^2)\langle\beta,\mathrm{out}|\phi(x)|\alpha,\mathrm{in}\rangle.
 \label{eq:oneleg}
\end{align}
The first term annihilates one incoming particle. Its packet overlaps are
the unchanged-particle or spectator contributions. The second term is the
reduced interacting part with one external particle replaced by a local
field insertion.

For an incoming packet, Hermitian conjugation of \eqref{eq:adifference}
and rearrangement give
\begin{equation}
 a_{\rm in}^\dagger(f)=a_{\rm out}^\dagger(f)
 +\frac{\ii}{\sqrt{\Zp}}\int\dd^4x\,f(x)(\Box+M^2)\phi(x).
\end{equation}
Both incoming and outgoing reductions therefore contribute \(+\ii/\sqrt{\Zp}\).
The difference is the wave: \(e^{-\ii k\cdot x}\) for an incoming particle
and \(e^{+\ii p\cdot x}\) for an outgoing one.

Repeated reduction orders the late annihilation operators to the left
and the early creation operators to the right. Interpolating between these
time boundaries produces a time-ordered vacuum correlator. Derivatives of
its step functions must be kept. For example, with a canonically normalized
scalar at the regulator level,
\begin{align}
 \partial_{x^0}^2\Tord\phi(x)\phi(y)
 &=\Tord\partial_{x^0}^2\phi(x)\phi(y)
 +\delta(x^0-y^0)[\dot\phi(x),\phi(y)],\\
 [\dot\phi(t,\bm x),\phi(t,\bm y)]&=-\ii\delta^{(3)}(\bm x-\bm y).
\end{align}
In the free case this yields
\((\Box_x+M^2)D_F(x-y)=-\ii\delta^{(4)}(x-y)\), even though the
operator field solves the homogeneous equation. A rescaled field has the
correspondingly rescaled contact coefficient.

\begin{keyidea}[Do not discard the identity process]
The overlap terms generated during reduction reconstruct unchanged
particles and products of independent processes. We use \(G_{N,c}\) to
extract the connected collision amplitude and add the identity and other
disconnected S-matrix contributions separately. In a free Gaussian theory
\(G_{N,c}=0\) for \(N>2\), so the collision amplitudes vanish while \(S=1\).
\end{keyidea}

\module{The scalar LSZ formula in position and momentum space}
Let \(k_1,\ldots,k_r\) be incoming momenta and \(p_1,\ldots,p_s\) outgoing
momenta, all future directed with mass \(M\), and set \(N=r+s\).
Iteration of the preceding reduction gives
\begin{align}
 \langle p_1\cdots p_s|S-1|k_1\cdots k_r\rangle_c
 &=\left(\frac{\ii}{\sqrt{\Zp}}\right)^N
 \int\prod_{a=1}^s\dd^4x_a\prod_{b=1}^r\dd^4y_b\,
 e^{\ii\sum_a p_a x_a-\ii\sum_b k_b y_b}\notag\\
 &\quad\times\prod_a(\Box_{x_a}+M^2)
 \prod_b(\Box_{y_b}+M^2)\,
 G_{N,c}(x_1,\ldots,x_s,y_1,\ldots,y_r).
 \label{eq:positionLSZ}
\end{align}
The plane-wave expression is shorthand for the packet version followed by
its distributional limit. Distinct stable species use their own masses,
fields, and residues on each external leg.

To avoid a momentum-sign ambiguity, define Fourier labels explicitly:
\begin{equation}
 \begin{gathered}
 \widehat G_{N,c}(q_1,\ldots,q_N)
 =\int\prod_j\dd^4z_j\,e^{\ii\sum_j q_j z_j}G_{N,c}(z_1,\ldots,z_N),\\
 (q_1,\ldots,q_N)=(p_1,\ldots,p_s,-k_1,\ldots,-k_r).
 \end{gathered}
\end{equation}
These are \emph{all-outgoing} Fourier labels; the physical incoming
momenta \(k_b\) still have positive energies. Integration by parts replaces
\(\Box+M^2\) by \(-(q^2-M^2)\), hence
\begin{equation}
 \boxed{\quad
 \ii(2\pi)^4\delta^{(4)}\!\left(\sum_jq_j\right)\cM
 =\left.\prod_{j=1}^N\frac{q_j^2-M^2}{\ii\sqrt{\Zp}}\,
 \widehat G_{N,c}(q_1,\ldots,q_N)\right|_{\text{external poles}}.\quad}
 \label{eq:momentumLSZ}
\end{equation}
One first extracts the pole coefficient, then evaluates it on the chosen
physical energy branches. Terms lacking at least one of the simultaneous
external poles do not contribute to this residue.

Equivalently, the leading multiple-pole part has the form
\begin{equation}
 \widehat G_{N,c}\sim (2\pi)^4\delta^{(4)}\!\left(\sum_jq_j\right)
 \left[\prod_j\frac{\ii\sqrt{\Zp}}{q_j^2-M^2+\ii0}\right]\ii\cM.
 \label{eq:factorization}
\end{equation}
Each local insertion overlaps with the external particle by \(\sqrt{\Zp}\).
This is why there is one square root per leg rather than one full residue
per leg in \eqref{eq:momentumLSZ}.

\begin{remarkbox}[The inverse operator is not the propagator]
\(\Box+M^2\) is the Klein--Gordon kinetic operator. Its Feynman inverse is
related to \(D_F\) by \((\Box+M^2)D_F=-\ii\delta\).
Calling both objects a ``free propagator'' hides precisely the inverse
operation that reduction performs.
\end{remarkbox}

\module{Amputated kernels, residue factors, and field choices}
Remove the overall momentum delta function and call the remainder of the
connected correlator \(\mathcal G_{N,c}\). Define a full-propagator-amputated
connected kernel \(\ii\mathcal V_N\) by
\begin{equation}
 \mathcal G_{N,c}(q_1,\ldots,q_N)
 =\left[\prod_j\widetilde G_2(q_j)\right]\ii\mathcal V_N(q_1,\ldots,q_N).
\end{equation}
Use the pole form \eqref{eq:pole} in \eqref{eq:momentumLSZ}. Each leg gives
\((q^2-M^2)\widetilde G_2/(\ii\sqrt{\Zp})\to\sqrt{\Zp}\), so
\begin{equation}
 \boxed{\quad\cM=\Zp^{N/2}\,\mathcal V_N\big|_{\rm on\ shell}.\quad}
 \label{eq:kernelLSZ}
\end{equation}
For four identical scalar external legs this is \(\cM=\Zp^2\mathcal V_4\),
consistent with Lecture 22. If the renormalized field is chosen to have
unit pole residue, these external factors are one.

Amputation with full propagators already removes their entire chains of
self-energy insertions. Adding those same chains again as external-leg
corrections double counts them. Internal propagators and subdiagrams are
still part of \(\mathcal V_N\). The connected amputated kernel is generally
\emph{not} identical to a 1PI vertex: a tree with two interaction vertices
joined by an internal line is connected but one-particle reducible.

A simple field rescaling checks the normalization. If \(\phi'=c\phi\)
with \(c>0\), then
\begin{equation}
 G'_{N,c}=c^NG_{N,c},\qquad Z'_{\rm pole}=c^2\Zp.
\end{equation}
Every factor of \(c\) cancels in \eqref{eq:momentumLSZ}. Negative or complex
phase choices merely require a consistent phase convention for the
particle states and field overlaps; observables remain unchanged.

More generally, take a local operator
\begin{equation}
 O=(1+\alpha\Box)\phi,\qquad [\alpha]=({\rm mass})^{-2}.
\end{equation}
Its one-particle overlap is \((1-\alpha M^2)\sqrt{\Zp}\). It interpolates
the same particle if that coefficient is nonzero, with the corresponding
residue and phase. Derivatives in time-ordered correlators can generate
contact terms, but they do not alter the simultaneous external particle
pole coefficient. If \(\alpha=M^{-2}\), this particular operator loses the
pole and cannot serve as an interpolating field for that particle.

\begin{keyidea}[Particles need not be elementary fields]
A renormalized composite operator with a nonzero stable-particle overlap
can also be used in reduction. The asymptotic particle and its amplitude
are physical; the interpolating operator is a choice of probe. Composite
operator renormalization and contact terms must be handled consistently.
\end{keyidea}

\module{A direct tree-level test, followed by the loop prescription}
Use the interaction \(-\lambda\phi^4/4!\) and its leading connected
four-point function from Lecture 21:
\begin{equation}
 G_{4,c}^{(1)}(x_1,x_2,x_3,x_4)
 =-\ii\lambda\int\dd^4z\,\prod_{j=1}^4D_F(x_j-z).
 \label{eq:treeG}
\end{equation}
At this order \(\Zp=1\), and the free mass agrees with the physical mass
to the accuracy needed for the tree amplitude. Apply the four operators
\(K_j=\Box_{x_j}+M^2\):
\begin{equation}
 \prod_jK_jG_{4,c}^{(1)}
 =(-\ii\lambda)(-\ii)^4\int\dd^4z\,
 \prod_j\delta^{(4)}(x_j-z).
\end{equation}
The four LSZ factors contribute \(\ii^4\). For two incoming and two
outgoing waves, the remaining integral is
\begin{equation}
 -\ii\lambda\int\dd^4z\,e^{\ii(p_3+p_4-k_1-k_2)\cdot z}
 =-\ii\lambda(2\pi)^4\delta^{(4)}(p_3+p_4-k_1-k_2).
\end{equation}
Thus \(\ii\cM=-\ii\lambda\), with no extra external \(2E\) factors and no
external propagators. The \(2E\)'s instead enter the state measure used
when converting this amplitude to a cross section.

At loop order the logic is unchanged, but each ingredient must be
calculated consistently:
\begin{enumerate}
 \item Locate the stable pole using \(M^2-m^2-\Sigma_R(M^2)=0\).
 \item Determine \(\Zp=[1-\Sigma_R'(M^2)]^{-1}\) for the chosen field.
 \item Compute the connected amputated kernel, including internal loops
 and counterterms to the requested order.
 \item Multiply by the external square roots of the residues and evaluate
 the external momenta on their physical mass shells.
\end{enumerate}
For massive \(\lambda\phi^4\) in the symmetric vacuum, the one-loop tadpole
is momentum independent. It shifts the mass but produces no derivative
\(\Sigma_R'\) at that order in a scheme with no finite field rescaling.
Hence \(\Zp=1+O(\lambda^2)\); its effect on the tree four-point amplitude
first appears at \(O(\lambda^3)\), beyond the one-loop four-point result.

\begin{remarkbox}[Vacuum normalization is a separate operation]
The ratio \(Z[J]/Z[0]\) removes vacuum components, taking connected
functions removes disconnected external components, and LSZ removes
external propagation and field overlaps. These operations solve different
problems. None is a substitute for the other two.
\end{remarkbox}

\module{External fermions: spinor overlaps and ordered amplitudes}
The particle-pole argument extends to spin. For a stable Dirac particle
of mass \(m\), choose conventions
\begin{align}
 \bra{\Omega}\psi_\alpha(0)\ket{p,s}&=\sqrt{Z_\psi}\,u_\alpha(p,s),\\
 \bra{\Omega}\bar\psi_\alpha(0)\ket{\bar p,s}
 &=\sqrt{Z_\psi}\,\bar v_\alpha(p,s),
\end{align}
where \(\ket{\bar p,s}\) denotes an antiparticle state of positive-energy
momentum \(p\). Near the pole,
\begin{equation}
 \widetilde S_F(p)\sim\frac{\ii Z_\psi(\slashed p+m)}{p^2-m^2+\ii0},\qquad
 \sum_su(p,s)\bar u(p,s)=\slashed p+m,\quad
 \sum_sv(p,s)\bar v(p,s)=\slashed p-m.
\end{equation}
Amputate the full external fermion propagators before imposing the
on-shell Dirac equation; then attach the appropriate spinor and one
\(\sqrt{Z_\psi}\) per external leg to the resulting kernel.
At tree level, or in a unit-pole-residue scheme, those square roots are one.

\begin{center}
\begin{tabular}{lll}\toprule
External state & Incoming & Outgoing\\\midrule
Scalar & \(1\) & \(1\)\\
Fermion & \(u(p,s)\) & \(\bar u(p,s)\)\\
Antifermion & \(\bar v(p,s)\) & \(v(p,s)\)\\\bottomrule
\end{tabular}
\end{center}
All entries refer to physical positive-energy momenta. Incoming spacetime
waves are \(e^{-\ii p\cdot x}\) and outgoing waves are \(e^{+\ii p\cdot x}\).
Grassmann differentiation and the chosen ordering of external Fock states
fix the signs. Exchanging two identical external fermions changes the
amplitude's sign; a closed fermion loop contributes its separate minus sign.

For example, in \(\Lag_{\rm int}=-y\phi\bar\psi\psi\), exchange of a scalar
between a fermion \(p_1\to p_3\) and another fermion \(p_2\to p_4\) gives
\begin{equation}
 \ii\cM_t=\big[\bar u(p_3)(-\ii y)u(p_1)\big]
 \frac{\ii}{(p_1-p_3)^2-m_\phi^2+\ii0}
 \big[\bar u(p_4)(-\ii y)u(p_2)\big].
\end{equation}
For identical fermions the full amplitude is this term minus the term
with \(p_3,s_3\leftrightarrow p_4,s_4\), for a fixed ordering of the states.
Assume stable external fermions; no external scalar is required here.

After constructing the amplitude, use the spin sums to compute
\(\overline{|\cM|^2}\), averaging over four incoming spin choices for two
unpolarized Dirac fermions. The flux and final-state phase-space formulas
of Lecture 24 then apply, including identical final-state counting.

\module{What reduction establishes, and what unitarity adds}
LSZ turns the singular part of a local correlator into an observable
amplitude. It does not say that every internal line is an on-shell particle.
Internal propagators carry momenta fixed by the diagram's conservation laws
and loop integrations; they are generally off shell. Real final particles
instead enter the phase-space sum with positive energies and their exact
mass shells.

Unitarity then constrains the amplitude extracted by reduction:
\begin{equation}
 -\ii\big(\cM_{fi}-\cM_{if}^*\big)
 =\sum_X\frac1{\cS_X}\int\dd\Phi_X\,
 \cM_{Xf}^*\cM_{Xi}.
 \label{eq:generalunitarity}
\end{equation}
Spin sums and the appropriate common total momentum are implicit.
For connected two-particle initial and final channels, this is the usual
amplitude form. In general multiparticle channels, retain the disconnected
and spectator pieces required by the operator identity. Only on the
diagonal \(f=i\) does the left side immediately become
\(2\operatorname{Im}\cM_{ii}\), giving the forward optical theorem.

When an amplitude is real analytic across a physical channel in the
appropriate sense, its boundary values obey
\begin{equation}
 \operatorname{Disc}_s\cM
 =\cM(s+\ii0)-\cM(s-\ii0)=2\ii\operatorname{Im}\cM(s+\ii0).
\end{equation}
For more general spin or channel amplitudes, Hermitian analyticity relates
the conjugate to the reversed process; it is not safe to assert the last
equality for every component without qualification. A discontinuity is
also not the same object as an imaginary part: they differ by \(2\ii\)
when that relation applies.

The same analytic correlator can yield different physical channels by
assigning different external energy branches and particle/antiparticle
labels. This motivates crossing. It is an analytic continuation between
physical regions, not a claim that a negative-energy particle is prepared
in the laboratory. Its general proof requires more than the elementary
pole reduction used here.

\begin{keyidea}[The completed chain of reasoning]
A stable mass shell and a nonzero field overlap identify an external
particle. Asymptotic packets define the S-matrix. Normalized path integrals
compute correlators, and LSZ extracts their physical pole amplitudes.
Phase space and flux turn those amplitudes into rates. Unitarity checks
that the resulting outcomes account for all probability.
\end{keyidea}
\textbf{Self-check.} Explain why changing the normalization of \(\phi\)
changes its correlators but not its scattering cross sections, why a
resonance cannot be reduced as an exact stable external state, and why
using the operator equation of motion inside a time-ordered product without
contact terms would fail even in the free theory.
\end{document}
